\SourceLeaf{049}
\noindent Perinde fit
\[
\begin{aligned}
x&=A+[\alpha\alpha]\xi^*+[\alpha\beta]\eta^*+[\alpha\gamma]\zeta^*+\text{etc.}\\
&\qquad-v^*(f[\alpha\alpha]+g[\alpha\beta]+h[\alpha\gamma]+\text{etc.})\\
&=A+[\alpha\alpha]\xi^*+[\alpha\beta]\eta^*+[\alpha\gamma]\zeta^*+\text{etc.}-Fv^*\\
&=A+[\alpha\alpha]\xi^*+[\alpha\beta]\eta^*+[\alpha\gamma]\zeta^*+\text{etc.}\\
&\qquad-\frac{F}{1+\omega}(F\xi^*+G\eta^*+H\zeta^*+\text{etc.}+K)
\end{aligned}
\]
Hinc itaque colligimus
\[
A^*=A-\frac{FK}{1+\omega}
\]
qui erit valor maxime plausibilis ipsius \(x\) ex \textit{omnibus} observationibus;
\[
[\alpha\alpha^*]=[\alpha\alpha]-\frac{FF}{1+\omega}
\]
adeoque pondus istius determinationis
\[
=\frac1{[\alpha\alpha]-\dfrac{FF}{1+\omega}}
\]
Prorsus eodem modo invenitur valor maxime plausibilis ipsius \(y\), \textit{omnibus} observationibus superstructus
\[
B^*=B-\frac{GK}{1+\omega}
\]
atque pondus huius determinationis
\[
=\frac1{[\beta\beta]-\dfrac{GG}{1+\omega}}
\]
et sic porro. Q.\ E.\ I.

Liceat huic solutioni quasdam annotationes adiicere.

I. Substitutis his novis valoribus \(A^*,B^*,C^*\) etc., functio \(v^*\) obtinet valorem maxime plausibilem
\[
K-\frac{K}{1+\omega}(Ff+Gg+Hh+\text{etc.})=\frac{K}{1+\omega}
\]
Et quum indefinite sit
\[
v^*=\frac{F}{1+\omega}.\xi^*+\frac{G}{1+\omega}.\eta^*+\frac{H}{1+\omega}.\zeta^*+\text{etc.}+\frac{K}{1+\omega}
\]
\SourceLeaf{050}%
pondus istius determinationis per principia art.\ 29 eruitur
\[
=\frac{1+\omega}{Ff+Gg+Hh+\text{etc.}}=\frac1\omega+1
\]
Eadem immediate resultant ex applicatione regulae in fine art.\ 21 traditae; scilicet complexus aequationum primitivarum praebuerat determinationem \(v^*=K\) cum pondere \(=\frac1\omega\), dein observatio nova dedit determinationem aliam, ab illa independentem, \(v^*=0\), cum pondere \(=1\), quibus combinatis prodit determinatio \(v^*=\frac{K}{1+\omega}\) cum pondere \(=\frac1\omega+1\).

II. Hinc porro sequitur, quum pro \(x=A^*,\allowbreak y=B^*,\allowbreak z=C^*\) etc. esse debeat \(\xi^*=0,\allowbreak\eta^*=0,\allowbreak\zeta^*=0\) etc., pro iisdem valoribus fieri
\[
\xi=-\frac{fK}{1+\omega},\quad\eta=-\frac{gK}{1+\omega},\quad\zeta=-\frac{hK}{1+\omega}\quad\text{etc.}
\]
nec non, quoniam indefinite \(\Omega=\xi(x-A)+\eta(y-B)+\zeta(z-C)+\text{etc.}+M\),
\[
\begin{aligned}
\Omega&=\frac{KK}{(1+\omega)^2}(Ff+Gg+Hh+\text{etc.})+M\\
&=M+\frac{\omega KK}{(1+\omega)^2}
\end{aligned}
\]
denique, quoniam indefinite \(\Omega^*=\Omega+v^*v^*\),
\[
\Omega^*=M+\frac{\omega KK}{(1+\omega)^2}+\frac{KK}{(1+\omega)^2}=M+\frac{KK}{1+\omega}
\]
III. Comparando haec cum iis, quae in art.\ 30 docuimus, animadvertimus, functionem \(\Omega\) hic valorem minimum obtinere, quem pro valore determinato functionis \(v^*=\frac{K}{1+\omega}\) accipere potest.

\Article{36}
Problematis alius, praecedenti affinis, puta

\textit{Investigare mutationes valorum maxime plausibilium incognitarum, a mutato pondere unius ex observationibus primitivis oriundas, nec non pondera novarum determinationum}

\noindent solutionem tantummodo hic adscribemus, demonstrationem, quae ad instar art. praec. facile absolvitur, brevitatis caussa supprimentes.

Supponamus, peracto demum calculo animadverti, alicui observationum pondus seu nimis parvum, seu nimis magnum tributum esse, e.\ g. observationi primae, quae dedit \(V=L\), loco ponderis \(p\) in calculo adhibiti rectius tribui pondus \(=p^*\). Tunc haud opus erit calculum integrum repetere, sed commodius correctiones per formulas sequentes computare licebit.

\SourceLeaf{051}%
Valores incognitarum maxime plausibiles correcti erunt hi:
\[
\begin{aligned}
x&=A-\frac{(p^*-p)\alpha\lambda}{p+(p^*-p)(a\alpha+b\beta+c\gamma+\text{etc.})}\\[7pt]
y&=B-\frac{(p^*-p)\beta\lambda}{p+(p^*-p)(a\alpha+b\beta+c\gamma+\text{etc.})}\\[7pt]
z&=C-\frac{(p^*-p)\gamma\lambda}{p+(p^*-p)(a\alpha+b\beta+c\gamma+\text{etc.})}
\end{aligned}
\]
etc. ponderaque harum determinationum invenientur, dividendo unitatem resp. per
\[
\begin{aligned}
&[\alpha\alpha]-\frac{(p^*-p)\alpha\alpha}{p+(p^*-p)(a\alpha+b\beta+c\gamma+\text{etc.})}\\[7pt]
&[\beta\beta]-\frac{(p^*-p)\beta\beta}{p+(p^*-p)(a\alpha+b\beta+c\gamma+\text{etc.})}\\[7pt]
&[\gamma\gamma]-\frac{(p^*-p)\gamma\gamma}{p+(p^*-p)(a\alpha+b\beta+c\gamma+\text{etc.})}\quad\text{etc.}
\end{aligned}
\]
Haec solutio simul complectitur casum, ubi peracto calculo percipitur, unam ex observationibus omnino reiici debuisse, quum hoc idem sit ac si facias \(p^*=0\); et perinde valor \(p^*=\infty\) refertur ad casum eum, ubi aequatio \(V=L\), quae in calculo tamquam approximata tractata erat, revera praecisione absoluta gaudet.

Ceterum quoties vel aequationibus, quibus calculus superstructus erat, \textit{plures} novae accedant, vel \textit{pluribus} ex illis pondera erronea tributa esse percipitur, computus correctionum nimis complicatus evaderet; quocirca in tali casu calculum ab integro reficere praestabit.

\Article{37}
In artt.\ 15.\ 16 methodum explicavimus, observationum praecisionem proxime determinandi\footnote{Disquisitio de eodem argumento, quam in commentatione anteriore (\textit{Bestimmung der Genauigkeit der Beobachtungen. Zeitschrift für Astronomie und verwandte Wissenschaften} Vol.\ I, p.\ 185) tradideramus, eidem hypothesi circa indolem functionis probabilitatem errorum exprimentis innixa erat, cui in Theoria motus corporum coelestium methodum quadratorum minimorum superstruxeramus (vid. art.\ 9, III).}. Sed haec methodus supponit, errores, qui revera occurrerint, satis multos exacte cognitos esse, quae conditio, stricte loquendo, rarissime, ne dicam numquam, locum habebit. Quodsi quidem quantitates, quarum valores approximati per observationes innotuerunt, secundum legem cognitam, ab una pluribusve quantitatibus incognitis pendent, harum valores maxime plausibiles per methodum quadratorum minimorum eruere licebit, ac dein valores quantitatum, quae observationum obiecta fuerant, illinc computati perparum a valoribus
